\chapter{Core Modules}
\label{ch:core_modules}

\section{Introduction}

This chapter documents the core modules that implement orbital mechanics algorithms. Each module is designed to be independent yet composable.

\subsection{KeplerianStateTyped}

The primary way to represent orbital elements in AstDyn is via the \texttt{KeplerianStateTyped<Frame>} template. This ensures that the elements are always bound to a specific reference frame and use type-safe units.

\begin{lstlisting}[language=C++,caption={KeplerianStateTyped class usage}]
namespace astdyn::physics {

// Example: Creating elements from a database (AU, deg) in Ecliptic J2000
auto elem = KeplerianStateTyped<core::ECLIPJ2000>::from_traditional(
    time::EpochTDB::from_mjd(60000.5),
    2.77,                // semi-major axis [AU]
    0.075,               // eccentricity
    10.6,                // inclination [deg]
    80.3,                // Omega [deg]
    73.6,                // omega [deg]
    0.0,                 // Mean anomaly [deg]
    GravitationalParameter::sun()
);

// Derived quantities (type-safe)
Distance q = elem.perihelion_distance();
double period = elem.period_days();

// Convert to Cartesian (Frame is preserved automatically)
auto cart = keplerian_to_cartesian(elem);
\end{lstlisting}

\subsection{CartesianStateTyped}

Position and velocity are represented by \texttt{CartesianStateTyped<Frame>}, which uses \texttt{Distance} and \texttt{Velocity} types for its components.

\begin{lstlisting}[language=C++,caption={CartesianStateTyped usage}]
using namespace astdyn::physics;
using namespace astdyn::core;

// Create from SI components
auto state = CartesianStateTyped<GCRF>::from_si(
    epoch,
    rx, ry, rz, // meters
    vx, vy, vz  // m/s
);

// Access components with unit conversion
double dist_au = state.position.x().to_au();
double speed_kms = state.velocity.norm().to_km_s();

// Explicit frame transformation (Compile-time verified)
auto ecliptic_state = state.cast_frame<ECLIPJ2000>();
\end{lstlisting}

\section{Force Models}

\subsection{ForceModel Interface}

\begin{lstlisting}[language=C++]
class ForceModel {
public:
    virtual ~ForceModel() = default;
    
    /**
     * @brief Compute acceleration [AU/day^2]
     * @param t Time (TDB)
     * @param pos_au Particle position [AU]
     * @param vel_au_d Particle velocity [AU/day]
     * @return Acceleration [AU/day^2]
     */
    virtual Eigen::Vector3d compute_acceleration(
        time::EpochTDB t, 
        const Eigen::Vector3d& pos_au, 
        const Eigen::Vector3d& vel_au_d) const = 0;
};
\end{lstlisting}

\subsection{Point Mass Gravity}

\begin{lstlisting}[language=C++]
class PointMassGravity : public ForceModel {
private:
    std::shared_ptr<ephemeris::PlanetaryEphemeris> ephem_;
    
public:
    PointMassGravity(std::shared_ptr<ephemeris::PlanetaryEphemeris> eph)
        : ephem_(eph) {}
    
    Eigen::Vector3d compute_acceleration(
        time::EpochTDB t, 
        const Eigen::Vector3d& pos_au, 
        const Eigen::Vector3d& vel_au_d) const override {
        
        // Multi-body planetary perturbation logic
        return ephem_->compute_acceleration(t, pos_au);
    }
};
\end{lstlisting}

\subsection{ForceField Aggregator}

The \texttt{ForceField} class manages multiple force models and provides the total acceleration required by the integrator.

\begin{lstlisting}[language=C++]
class ForceField {
public:
    Eigen::Vector3d total_acceleration(
        time::EpochTDB t, 
        const Eigen::Vector3d& pos_au, 
        const Eigen::Vector3d& vel_au_d) const;
};
\end{lstlisting}

\section{Numerical Integration}

\subsection{Integrator Interface}

\begin{lstlisting}[language=C++]
class Integrator {
public:
    virtual ~Integrator() = default;
    
    virtual Eigen::VectorXd integrate(
        const DerivativeFunction& f,
        const Eigen::VectorXd& y0,
        double t0,
        double tf) = 0;
};
\end{lstlisting}

Standard implementations include \texttt{RK4Integrator} (fixed step) and \texttt{RKF78Integrator} (adaptive step).

\section{Orbit Propagation}

The \texttt{Propagator} class orchestrates the force field and the integrator.

\begin{lstlisting}[language=C++]
class Propagator {
public:
    template <typename Frame>
    physics::KeplerianStateTyped<Frame> propagate_keplerian(
        const physics::KeplerianStateTyped<Frame>& state,
        time::EpochTDB target_time);

    template <typename Frame>
    physics::CartesianStateTyped<Frame> propagate_cartesian(
        const physics::CartesianStateTyped<Frame>& state,
        time::EpochTDB target_time);
};
\end{lstlisting}

\textbf{Usage Example}:
\begin{lstlisting}[language=C++]
// 1. Setup
auto ephem = std::make_shared<ephemeris::PlanetaryEphemeris>();
auto integrator = std::make_shared<RKF78Integrator>(0.1, 1e-12);

PropagatorSettings settings;
settings.include_planets = {"Sun", "Jupiter", "Saturn"};

Propagator prop(integrator, ephem, settings);

// 2. Propagate
auto final_state = prop.propagate_keplerian(initial_state, t_final);
\end{lstlisting}

\textbf{Usage Example}:
\begin{lstlisting}[language=C++]
// Setup
auto spice = std::make_shared<SpiceInterface>();
spice->load_kernel("de440.bsp");

auto forces = std::make_shared<PointMassGravity>(
    spice, {Body::SUN, Body::JUPITER, Body::SATURN});

auto integrator = std::make_shared<RKF78>(1e-12);

Propagator prop(integrator, forces, spice);

// Propagate Pompeja for 60 days
Vector6d y0 = /* initial state */;
double t0 = 2460000.5;
double t1 = t0 + 60.0;

Vector6d y_final = prop.propagate(y0, t0, t1);

std::cout << "Final position: " << y_final.head<3>().transpose() << " AU\n";
\end{lstlisting}

\section{Observations}

\subsection{OpticalObservation Class}

\begin{lstlisting}[language=C++]
namespace astdyn::observations {

struct OpticalObservation {
    std::string object_designation;
    std::string observatory_code;
    
    time::EpochUTC time;             ///< Time of observation (UTC)
    
    astrometry::RightAscension ra;   ///< Right Ascension
    astrometry::Declination dec;    ///< Declination
    astrometry::Angle sigma_ra;      ///< RA uncertainty
    astrometry::Angle sigma_dec;     ///< Dec uncertainty
    
    std::optional<double> magnitude; ///< Apparent magnitude
};

} // namespace
\end{lstlisting}

\subsection{MPC Reader}

The \texttt{MPCReader} parses Minor Planet Center 80-column and ADES formats.

\begin{lstlisting}[language=C++]
auto observations = io::MPCReader::read_file("observations.txt");
\end{lstlisting}

\section{Observatory Database}

The \texttt{ObservatoryDatabase} manages locations using MPC codes.

\begin{lstlisting}[language=C++]
struct Observatory {
    std::string code;
    double longitude;  ///< East longitude [rad]
    double latitude;   ///< Latitude [rad]
    double altitude;   ///< Altitude [m]
    
    // Geocentric position at given time
    math::Vector3<core::GCRF, physics::Distance> getPositionGCRF(time::EpochUTC t) const;
};
\end{lstlisting}

\section{Summary}

Core modules provide:

\begin{enumerate}
    \item \textbf{Orbital Elements}: Keplerian, Cartesian, Cometary representations
    \item \textbf{Force Models}: Extensible interface for perturbations
    \item \textbf{Integrators}: Adaptive step-size RK methods
    \item \textbf{Propagator}: High-level orbit propagation with STM
    \item \textbf{Observations}: Astrometric measurements and MPC parsing
    \item \textbf{Observatories}: Geodetic coordinates and transformations
\end{enumerate}

All modules are designed for composition and extensibility.
